% page 547 of the facsimile (image p-546 of the scan)
% block 170.2 x 250.6 mm ; left margin 31.8 mm
\pgc{10.160mm}{0.000mm}{
\l{\cel{51}539}
\l{smalto.\cel{1}\sou{(tek}n.)\cel{2}Vitro blua, obtenata per la fuzo di materio}
\l{\cel{3}vitrigebla kun kobalto. - DEFIRS.}
\l{}
\l{smeraldo. Lapido diafana, verda. - DEFIRS.}
\l{}
\l{}
\l{smerilo. Petro harda, ek alumino, silico e fer-oxido, pulverigita}
\l{\cel{3}por polisar la lapidi, la petri, la metali, kristalo. - FIS.}
\l{}
\l{\sou{snapa}r. (trans.) Sizar bruske, per un maxil-ago. DE.}
\l{}
\l{sniflar. (trans.) Aspirar ulo per sua nazo (tabako pulverigita,}
\l{\cel{3}aquo, la humoro qua stagnas en la naztrui). - E.}
\l{}
\l{snortar. (netrans.) (Dicesas precipue pri kavalo). Sniflar tre}
\l{\cel{3}bruisoze por manifestar sua repugneso ; ekpulsar aero per sua}
\l{\cel{3}naztrui tre bruisoze por manifestar sua pavoro, nepacienteso,}
\l{\cel{3}e c. - E.}
\l{}
\l{\sou{sobr}a. I. Qua su moderas pri manjo e drinko. - II.\cel{1}\sou{(met}af.)}
\l{\cel{3}Qua esas moderata pri sua\cel{1}\sou{agi}, ne exajerema. (Ex.: Lu esas}
\l{\cel{3}sobra pri komplimentar). - EFIS.}
\l{}
\l{socio. Ensemblo,duroza, ek homi qui vivas obedie di legi komuna.}
\l{\cel{3}\sou{(an}ke aludante ula sorti de animali qui vivas ensemble : abeli,}
\l{\cel{3}formiki, e c.) - DEFIRS.}
\l{}
\l{sociologo. Ciencisto qua su okupas pri sociologio.\cel{2}DEFIRS}
\l{}
\l{sociologio. Cienco di qua la studiajo esas la cirkonstanci ed}
\l{\cel{3}existo-kondicioni di la socio homara. - DEFIRS.}
\l{}
\l{sociso. Peco de la intestino di porko, di mutono, quan onu plen-}
\l{\cel{3}igis ye karno de porko, o de bovo, tre spicizita, triturita e}
\l{\cel{3}pistita, e tre presita. - EFIR.}
\l{}
\l{sodo. (bot.)\cel{3}Planto ek la\cel{1}\sou{famil}io \textquotedbl{}kenopodiacei\textquotedbl{}, propra ye la}
\l{\cel{3}regioni klimato-moderata e subtropika ; un-yar-dura ; di qua}
\l{\cel{3}la folii esas quaze pulpoza, lineala, e la flori mikra e ver-}
\l{\cel{3}datra ; de qui onu extraktis olim natroxo. - L. salsola soda}
\l{\cel{3}e salsola kali. - (kemio) Salo alkalioza quan onu extraktas}
\l{\cel{3}de ta planto ed anke ek la fuki.\cel{1}\sou{Simbolo kemial}a : CO3 Na2 .}
\l{}
\l{\cel{3}DEFIR.}
\l{}
\l{sodomiar.\cel{1}\sou{(netrans}.) Agar koito anusala, inter-vira, od inter}
\l{\cel{3}viro e muliero, o kun bestio. - DEFIS.}
\l{}
\l{sofao. Sorto di repozo-lito, kun tri dors-apogili, uzata kom}
\l{\cel{3}sidilo.- DEFIS.}
\l{}
\l{sofismo. Argumento trompera, seduktera. - DEFIRS.}
\l{}
\l{sofisto. I. (en la\cel{1}\sou{epoki antiq}ua, en Grekia). Sorto de profesoro}
\l{\cel{3}di filozofio, di retoriko, qua instruktis pri quale pledar \textquotedbl{}por\textquotedbl{}}
\l{\cel{3}e \textquotedbl{}kontre\textquotedbl{} alterne, un tezo, sen egardo pri olua exakteso.-}
\l{\cel{3}II. (nun) Persono qua argumentas trompere, seduktera. - DEFIRS.}
}
